\begin{longtable}{|p{3.2cm}|p{4.2cm}|p{5.4cm}|p{2.5cm}|}
\caption{Types of Product Development Processes with Their Features and Examples}
\label{tab:session2-process-types} \\
\hline
\textbf{Process Type} & \textbf{Description} & \textbf{Distinct Features} & \textbf{Examples} \\
\hline
Generic (Market-Pull) Products & The team begins with a market opportunity and selects appropriate technologies to meet customer needs. & Process generally includes distinct planning, concept development, system-level design, detail design, testing and refinement, and production ramp-up phases. & Sporting goods, furniture, tools. \\
\hline
Technology-Push Products & The team begins with a new technology, then finds an appropriate market. & Planning phase involves matching technology and market. Concept development assumes a given technology. & Gore-Tex rainwear, Tyvek envelopes. \\
\hline
Platform Products & The team assumes that the new product will be built around an established technological subsystem. & Concept development assumes a proven technology platform. & Consumer electronics, computers, printers. \\
\hline
Process-Intensive Products & Characteristics of the product are highly constrained by the production process. & Either an existing production process must be specified from the start, or both product and process must be developed together from the start. & Snack foods, breakfast cereals, chemicals, semiconductors. \\
\hline
Customized Products & New products are slight variations of existing configurations. & Similarity of projects allows for a streamlined and highly structured development process. & Motors, switches, batteries, containers. \\
\hline
High-Risk Products & Technical or market uncertainties create high risks of failure. & Risks are identified early and tracked throughout the process. Analysis and testing activities take place as early as possible. & Pharmaceuticals, space systems. \\
\hline
Quick-Build Products & Rapid modeling and prototyping enables many design-build-test cycles. & Detail design and testing phases are repeated a number of times until the product is completed or time/budget runs out. & Software, cellular phones. \\
\hline
Product-Service Systems & Products and their associated service elements are developed simultaneously. & Both physical and operational elements are developed, with particular attention to design of the customer experience and the process flow. & Restaurants, software applications, financial services. \\
\hline
Complex Systems & System must be decomposed into several subsystems and many components. & Subsystems and components are developed by many teams working in parallel, followed by system integration and validation. & Airplanes, jet engines, automobiles. \\
\hline
\end{longtable}